\documentclass[11pt]{amsart}

\usepackage[T1]{fontenc}
\usepackage{lmodern}
\usepackage{mathtools,amssymb}
\usepackage{booktabs}
\usepackage[margin=0.85in]{geometry}
\usepackage{microtype}
\usepackage[hidelinks]{hyperref}

\newcommand{\Z}{\mathbb Z}
\newcommand{\B}{\mathcal B_{41}}
\newcommand{\A}{\mathcal A_{41}}
\newcommand{\PhiMap}{\Phi_{41}}

\title[The complete fibre table at $n=41$]
{The Complete Fibre Table at $n=41$\\
for the Two-to-One Tunnell Map}

\author{Aakash Gurung}
\address{Department of Mathematics, The University of Alabama, Tuscaloosa, AL 35487, USA}
\email{agurung3@crimson.ua.edu}

\author{Kyungyong Lee}
\address{Department of Mathematics, The University of Alabama, Tuscaloosa, AL 35487, USA}
\email{klee94@ua.edu}

\date{September 2026}

\hypersetup{
  pdftitle={The Complete Fibre Table at n=41 for the Two-to-One Tunnell Map},
  pdfauthor={Aakash Gurung and Kyungyong Lee}
}

\begin{document}

\maketitle

\section*{The representation sets}

This note accompanies the manuscript \emph{Stable Matching and a Two-to-One
Map for Tunnell's Ternary Forms}.  We spell out the complete map at $n=41$ in
the original ternary coordinates.  Thus
\[
 \B=\{(x,y,z)\in\Z^3:2x^2+y^2+8z^2=41\},
 \qquad
 \A=\{(u,v,w)\in\Z^3:2u^2+v^2+32w^2=41\}.
\]
There are $32$ points in $\B$ and $16$ points in $\A$.  The tables below list
every target and its two preimages, so all points of both representation sets
occur.

\section*{The three rules used in the table}

\textbf{Even branch.}
For every $(u,v,w)\in\A$, the first preimage is
\[
 (u,v,2w)\longmapsto(u,v,w).
\]
This is equation~(4) of the paper and is the upper bijection in the proof of
Theorem~1.2.

\textbf{Direct odd branch.}
All eight direct odd points at $n=41$ lie in the third quarter-turn domain.
Equation~(10) gives
\[
 \Lambda_3(x,y,z)=\left(2z,y,-\frac{x}{4}\right).
\]
The domain test is Proposition~3.1, and the fact that this is a bijection onto
the direct target set is Theorem~3.2.

\textbf{Residual odd branch.}
The remaining eight odd points are paired by the unique stable matching of
Proposition~4.5.  Equation~(12) fixes the sign when the orbit matching is
lifted back to individual points.  Algorithm~7.1 computes the same matching,
its correctness is Theorem~7.3, and the complete $n=41$ residual run is
Example~8.1.  The four canonical residual pairs in original coordinates are
\[
\begin{aligned}
 (-2,-5,-1)&\longmapsto(4,3,0),&
 (-2,-5,1)&\longmapsto(4,-3,0),\\
 (-2,5,-1)&\longmapsto(0,-3,1),&
 (-2,5,1)&\longmapsto(0,-3,-1).
\end{aligned}
\]
Their negatives give the other four residual pairs.

\section*{How to read a fibre}

For a target \(Y\in\A\), its \emph{fibre} is
\[
 \PhiMap^{-1}(Y)=\{X\in\B:\PhiMap(X)=Y\}.
\]
The map is designed so that this set has two elements.  One has even third
coordinate and comes from equation~(4).  The other has odd third coordinate
and comes either from a quarter-turn or from the residual matching.  In the
two examples below, we first check that every displayed point satisfies its
quadratic equation.  We then determine the applicable branch and compute the
image.

\newpage
\section*{Worked example 1: a direct quarter-turn fibre}

Take the target
\[
 Y=(2,-1,1).
\]

\noindent\textbf{Step 1: Check that \(Y\) is a target point.}
The target equation is \(2u^2+v^2+32w^2=41\).  Substitution gives
\[
 2(2)^2+(-1)^2+32(1)^2=8+1+32=41,
\]
so \(Y\in\A\).

\medskip
\noindent\textbf{Step 2: Find the even preimage.}
Equation~(4) says that the even preimage of \((u,v,w)\) is \((u,v,2w)\).
Therefore
\[
 X_{\mathrm{even}}=(2,-1,2).
\]
It really is a source point because
\[
 2(2)^2+(-1)^2+8(2)^2=8+1+32=41.
\]
Its third coordinate is even, and equation~(4) gives
\[
 \PhiMap(2,-1,2)=(2,-1,1)=Y.
\]

\medskip
\noindent\textbf{Step 3: Check the proposed odd preimage.}
Consider
\[
 X_{\mathrm{odd}}=(-4,-1,1).
\]
It belongs to \(\B\), since
\[
 2(-4)^2+(-1)^2+8(1)^2=32+1+8=41.
\]
Its third coordinate is odd, so the even rule cannot apply.

\medskip
\noindent\textbf{Step 4: Determine the odd branch.}
For \(X_{\mathrm{odd}}=(x,y,z)=(-4,-1,1)\), the first two quarter-turn
tests give
\[
 x+2z+y=-4+2-1=-3,
 \qquad
 y-x+2z=-1+4+2=5.
\]
Neither number is divisible by \(8\).  The third test does hold because
\(4\mid x\).  Proposition~3.1 shows that the three tests are mutually
exclusive, so equation~(10) is the applicable rule.

\medskip
\noindent\textbf{Step 5: Apply the quarter-turn formula.}
Using Theorem~3.2 and equation~(10),
\[
\begin{aligned}
 \PhiMap(-4,-1,1)
   &=\Lambda_3(-4,-1,1)\\
   &=\left(2(1),-1,-\frac{-4}{4}\right)\\
   &=(2,-1,1)=Y.
\end{aligned}
\]

\medskip
\noindent\textbf{Step 6: Write the complete fibre.}
The even branch is a bijection, so \(Y\) has only one even preimage.
Theorem~3.2 makes the direct odd branch a bijection onto its image, so it has
only one direct odd preimage.  The assembly in Theorem~1.2 shows that there
are no further preimages.  Hence
\[
 \boxed{\PhiMap^{-1}(2,-1,1)
 =\{(2,-1,2),\,(-4,-1,1)\}.}
\]

\newpage
\section*{Worked example 2: a residual matching fibre}

Take the target
\[
 Y=(4,-3,0).
\]

\noindent\textbf{Step 1: Check that \(Y\) is a target point.}
We have
\[
 2(4)^2+(-3)^2+32(0)^2=32+9=41,
\]
so \(Y\in\A\).

\medskip
\noindent\textbf{Step 2: Find the even preimage.}
Equation~(4) gives
\[
 X_{\mathrm{even}}=(4,-3,0),
 \qquad
 \PhiMap(4,-3,0)=(4,-3,0)=Y.
\]
The same calculation \(2(4)^2+(-3)^2+8(0)^2=41\) checks that this point
belongs to \(\B\).

\medskip
\noindent\textbf{Step 3: Check the proposed odd preimage and its branch.}
Consider
\[
 X_{\mathrm{odd}}=(-2,-5,1).
\]
It is a source point because
\[
 2(-2)^2+(-5)^2+8(1)^2=8+25+8=41.
\]
The direct tests are
\[
\begin{aligned}
 x+2z+y&=-2+2-5=-5,\\
 y-x+2z&=-5+2+2=-1,\\
 x&=-2.
\end{aligned}
\]
The first two values are not divisible by \(8\), and the last is not
divisible by \(4\).  Thus none of the three quarter-turns applies, and this
point must enter the residual matching.

\medskip
\noindent\textbf{Step 4: Follow the residual orbit through the matching.}
In the lifted coordinates of Example~8.1, \(X_{\mathrm{odd}}\) is the
canonical representative
\[
 E=(-2,-5,2)
\]
of the source orbit \(s_2\).  Its preference order is
\[
 t_1,\ t_4,\ t_2,\ t_3.
\]
During Algorithm~7.1, \(s_2\) first proposes to \(t_1\), but \(t_1\) keeps
\(s_1\), whose edge has the smaller key.  It next proposes to \(t_4\), which
keeps \(s_3\) for the same reason.  Its third proposal, to \(t_2\), is
accepted.  This is the \(s_2\leftrightarrow t_2\) pair in Example~8.1.
Proposition~4.5 proves that the resulting stable matching is the unique one,
and Theorem~7.3 proves that the algorithm computes it.

\medskip
\noindent\textbf{Step 5: Recover the sign of the individual target.}
The canonical representative of \(t_2\) is
\[
 F=(-4,3,0).
\]
The accepted signed edge has \(\eta=-1\).  Since \(E\) itself was the
canonical source representative, \(\sigma=1\), and equation~(12) gives
\[
 E\longmapsto \sigma\eta F=(1)(-1)F=-F=(4,-3,0).
\]
Returning from the lifted coordinates to the original ones therefore gives
\[
 \PhiMap(-2,-5,1)=(4,-3,0)=Y.
\]

\medskip
\noindent\textbf{Step 6: Write the complete fibre.}
Equation~(4) supplies the unique even preimage.  Proposition~4.5 and the sign
rule~(12) supply the unique residual odd preimage.  Theorem~1.2 excludes any
additional preimages.  Therefore
\[
 \boxed{\PhiMap^{-1}(4,-3,0)
 =\{(4,-3,0),\,(-2,-5,1)\}.}
\]

\section*{A quick third example by negation}

The residual construction is equivariant under negation: if
\(\PhiMap(X)=Y\), then \(\PhiMap(-X)=-Y\).  Negating the odd pair in the
second example immediately gives
\[
 (2,5,-1)\longmapsto(-4,3,0).
\]
The even preimage of \((-4,3,0)\) is again \((-4,3,0)\).  Hence
\[
 \PhiMap^{-1}(-4,3,0)
 =\{(-4,3,0),\,(2,5,-1)\}.
\]

\section*{Direct fibres}

In each row, the middle point is the even preimage supplied by equation~(4),
and the last point is the odd preimage supplied by Theorem~3.2 and
equation~(10).

\begin{center}
\small
\renewcommand{\arraystretch}{1.22}
\begin{tabular}{@{}c c c c@{}}
\toprule
target $Y\in\A$ & even preimage & odd preimage & direct computation \\
\midrule
$(-2,-1,-1)$ & $(-2,-1,-2)$ & $(4,-1,-1)$
  & $\Lambda_3(4,-1,-1)=Y$ \\
$(-2,-1, 1)$ & $(-2,-1, 2)$ & $(-4,-1,-1)$
  & $\Lambda_3(-4,-1,-1)=Y$ \\
$(-2, 1,-1)$ & $(-2, 1,-2)$ & $(4, 1,-1)$
  & $\Lambda_3(4,1,-1)=Y$ \\
$(-2, 1, 1)$ & $(-2, 1, 2)$ & $(-4,1,-1)$
  & $\Lambda_3(-4,1,-1)=Y$ \\
$( 2,-1,-1)$ & $( 2,-1,-2)$ & $(4,-1,1)$
  & $\Lambda_3(4,-1,1)=Y$ \\
$( 2,-1, 1)$ & $( 2,-1, 2)$ & $(-4,-1,1)$
  & $\Lambda_3(-4,-1,1)=Y$ \\
$( 2, 1,-1)$ & $( 2, 1,-2)$ & $(4,1,1)$
  & $\Lambda_3(4,1,1)=Y$ \\
$( 2, 1, 1)$ & $( 2, 1, 2)$ & $(-4,1,1)$
  & $\Lambda_3(-4,1,1)=Y$ \\
\bottomrule
\end{tabular}
\end{center}

The sixth row is the fibre computed step by step in Worked example~1.  The
other seven rows use the same two formulas.

\newpage
\section*{Residual fibres}

The notation $s_i$ refers to the four canonical source orbits in
Example~8.1.  A minus sign indicates that equation~(12) applies the
antipodal lift to the corresponding canonical pair.

\begin{center}
\small
\renewcommand{\arraystretch}{1.22}
\begin{tabular}{@{}c c c c@{}}
\toprule
target $Y\in\A$ & even preimage & odd preimage & residual origin \\
\midrule
$(-4,-3,0)$ & $(-4,-3,0)$ & $(2,5,1)$
  & negative of the $s_1$ pair \\
$(-4, 3,0)$ & $(-4, 3,0)$ & $(2,5,-1)$
  & negative of the $s_2$ pair \\
$(0,-3,-1)$ & $(0,-3,-2)$ & $(-2,5,1)$
  & $s_4$ pair \\
$(0,-3, 1)$ & $(0,-3, 2)$ & $(-2,5,-1)$
  & $s_3$ pair \\
$(0, 3,-1)$ & $(0, 3,-2)$ & $(2,-5,1)$
  & negative of the $s_3$ pair \\
$(0, 3, 1)$ & $(0, 3, 2)$ & $(2,-5,-1)$
  & negative of the $s_4$ pair \\
$(4,-3,0)$ & $(4,-3,0)$ & $(-2,-5,1)$
  & $s_2$ pair \\
$(4, 3,0)$ & $(4, 3,0)$ & $(-2,-5,-1)$
  & $s_1$ pair \\
\bottomrule
\end{tabular}
\end{center}

The stable matching in Example~8.1 is obtained from the proposal sequence
\[
\begin{aligned}
 s_1&\to t_1,&
 s_2&\to t_1\ \text{(rejected)},&
 s_3&\to t_4,\\
 s_4&\to t_3,&
 s_2&\to t_4\ \text{(rejected)},&
 s_2&\to t_2.
\end{aligned}
\]
It therefore matches the four source orbits as
\[
 s_1\leftrightarrow t_1,\qquad
 s_2\leftrightarrow t_2,\qquad
 s_3\leftrightarrow t_4,\qquad
 s_4\leftrightarrow t_3,
\]
\noindent
after which equation~(12) determines the individual signs displayed in the
table.

\section*{Exhaustiveness check}

The direct table contains the eight targets in the image of the quarter-turn
layer.  The residual table contains the other eight targets.  Their even
preimages are the $16$ points of $\B$ having even third coordinate; their odd
preimages consist of the eight direct points and the eight residual points.
Thus the tables contain $16$ distinct targets and $32$ distinct sources, and
each target has exactly two preimages, as asserted by Theorem~1.2.

The table was also regenerated from the formalized executable map
\texttt{tunnellMapTotal41}.  The Lean declarations
\texttt{card\_BRep\_41}, \texttt{card\_ARep\_41}, and
\texttt{tunnellMapTotal41\_exactly\_two} verify respectively that the source
and target cardinalities are $32$ and $16$ and that every displayed target
has exactly two preimages.

\end{document}
